\documentclass[conference]{IEEEtran}
\usepackage[utf8]{inputenc}
\usepackage[T1]{fontenc}
\usepackage{cite}
\usepackage{amsmath,amssymb}
\usepackage{graphicx}
\usepackage{booktabs}
\usepackage{url}

\title{Empirical Evaluation of a Deterministic-Validator Guard for LLM\ensuremath{\rightarrow}NetOps Generate\ensuremath{\rightarrow}Verify\ensuremath{\rightarrow}Repair Pipelines (Revised)}

\author{
\IEEEauthorblockN{Autonomous AI Researcher}
\IEEEauthorblockA{CNSM 2026 Agentic AI Researcher Track}
}

\begin{document}

\maketitle

\begin{abstract}
We preregistered and executed a paired confirmatory experiment that measures the operational impact of inserting a deterministic network-configuration validator into an LLM-driven generate\ensuremath{\rightarrow}verify\ensuremath{\rightarrow}repair (GVR) loop for intent\ensuremath{\rightarrow}configuration NetOps tasks. Using a reproducible synthetic 200-task multi-vendor intent bank and a single hosted model (gpt-5-mini) under an adapter contract (<=1 automated repair call per task), each task produced one shared initial candidate; the guarded artifact was the final configuration after deterministic validation and, if invoked, a single automated repair. Primary estimand: paired absolute difference in actionable-misconfiguration rate (guarded - baseline), implemented as paired success-rate difference per the preregistration. Results (N=200 pairs): baseline valid-config count = 199/200 (0.995), guarded valid-config count = 200/200 (1.000); absolute paired change (guarded - baseline) = 0.005 (0.5 percentage points). Exact McNemar two-sided test p = 1.0; paired bootstrap 95\% CI = [0.0, 0.015] (10,000 resamples, seed 1729). The preregistered primary hypothesis (>=50\% relative reduction in actionable misconfigurations under original power assumptions) is not supported because the observed baseline error prevalence (1/200) was far lower than assumed in the preregistered power calculation. We reconcile estimand algebra, correct contingency-cell notation, expand execution/accounting and reproducibility details, and point auditors to archived artifact paths and SHA-256 digests for forensic verification.
\end{abstract}

\section{Introduction}
Generative LLMs can translate operator intent into device-specific network configurations, promising reduced operator effort for routine NetOps tasks. However, incorrect or out-of-order commands can cause outages or violate workflow policies; production proposals therefore commonly interpose deterministic validators and automated repair loops as guardrails in generate\ensuremath{\rightarrow}verify\ensuremath{\rightarrow}repair (GVR) pipelines. Despite intuitive appeal, rigorous, preregistered quantification of the incremental operational benefit from adding a deterministic validator under a bounded adapter contract is scarce and often lacks forensic artifact traceability. We preregistered study cnsmspec2026\_gvr\_v1\_prereg\_001 (SHA-256 = 5cb8a30815eacf0a3ca659d1a54fbaa71218e5ce57c0b13e2658099eb7da6ee4) to answer: How much does integrating a deterministic network validator into an LLM-driven GVR pipeline reduce actionable misconfiguration rates (paired comparison) on a 200-task synthetic multi-vendor NetOps task bank, and what residual error types persist after automated repairs? Our primary methodological novelty is strictly forensic: (1) a preregistered paired estimand that compares, per task, the shared initial candidate against the final guarded artifact, (2) deterministic validator traces that emit canonical ASTs and simulator-backed semantic checks, and (3) cryptographic digests for all principal artifacts to permit deterministic auditing.

\section{Related Work}
Deterministic network verification and runtime validators (header-space analysis; Kazemian et al.; incremental checkers such as VeriFlow [VeriFlow DOI]) are well-established pre-deployment safety methods. Recent LLM\ensuremath{\rightarrow}NetOps evaluations (e.g., NetConfEval [NetConfEval DOI], Lira et al.) report generation accuracy and task-level metrics but typically omit preregistration, deterministic validator traces, or adapter-constrained repair budgets. Our work differs by binding the experiment to a single, archived adapter contract (hosted\_netops\_gvr\_v1), enforcing shared\_initial generation semantics (single initial candidate reused across conditions), and publishing forensic SHA-256 digests for execution and analysis artifacts so auditors can reconstruct every contingency cell and per-episode validator\_trace. We position our contribution as a narrow, artifact-grounded quantification of marginal validator benefit under a conservative repair budget rather than a broad model-comparison benchmark.

\section{Methodology}
Overview and preregistration. We implemented the preregistered paired binary design documented in cnsmspec2026\_gvr\_v1\_prereg\_001 (SHA-256 = 5cb8a30815eacf0a3ca659d1a54fbaa71218e5ce57c0b13e2658099eb7da6ee4). The preregistration specifies N=200 independent intent\ensuremath{\rightarrow}configuration tasks sampled from a deterministic synthetic task bank, paired observations per task (shared initial candidate used for both baseline and guarded conditions), and a single confirmatory estimand: paired\_success\_rate\_difference\_guarded\_minus\_baseline. The confirmatory test is an exact McNemar two‑sided test; a paired-bootstrap percentile 95\% CI (10,000 resamples; seed 1729) was prespecified to quantify uncertainty. The planned stratification (for sampling only) was Cisco-like N=66, Juniper-like N=67, RFC-neutral N=67 (see preregistration). All design elements (inclusion/exclusion rules, missingness policy, and multiplicity plan) were frozen prior to execution and are archived with SHA-256 digests in the execution manifest.

Adapter contract and generation semantics. Execution used adapter\_family = hosted\_netops\_gvr\_v1 with requested\_model = gpt-5-mini as declared in execution/model\_configuration.json (SHA-256 = a40e96e212ab87da251ac0d6dbb75bc543afa13438b5a6b9ebd073597ffdd820). Crucially, the adapter enforces shared\_initial\_generation semantics (independent\_condition\_generation = false): a single initial candidate was produced per task and that same candidate served as the baseline artifact and as the starting point for the guarded pipeline (so differences arise only where the guarded pipeline invoked the registered repair path). The adapter contract limited automated repair attempts to at most one per task (maximum\_repair\_calls\_per\_task = 1); provider-call accounting in deterministic\_reconciliation.json documents stage\_counts = \{"shared\_initial\_generation": 200, "repair": 1\} and hosted\_model\_call\_event\_count = 201 (see analysis/deterministic\_reconciliation.json SHA-256 = 8a2b85f8260409eebce1641421af6a74f1c479e505bdef805314510e25931891 and execution/raw\_results.jsonl SHA-256 = 9eeaa0c87db4742c807240b5ee18f934aaf45848350b3840f92af401bcb68d02). These enforced semantics remove cross-condition generation variance and isolate the deterministic validator+repair effect on the same initial LLM candidate.

Synthetic task bank generation and task encoding. The task bank was produced by deterministic\_netops\_task\_generator\_v5 (generator\_version = 5.0) and encodes for each task: initial\_state, expected\_final\_state, workflow policies (e.g., must\_be\_down\_for\_fields, group constraints), allowed\_touched\_fields, a human-readable intent, and an explicit required\_sequence of field-level operations. Task manifests are archived under execution/task\_manifest.jsonl (SHA-256 = 4b28215e8a611feda50f6284058083b9ffcc2f5881f557081ce9e9f890f1c78a). Representative task payloads in the archive show structured difficulty metadata (changed\_interface\_count, dependency\_rule\_count, pattern labels) which were used post hoc for exploratory stratified analyses; however the confirmatory inference is the paired binary test on the primary estimand described above.

Deterministic validator design and scoring semantics. The deterministic validator (deterministic\_netops\_validator\_v5, validator\_version = 5.0) implements a sequence of deterministic, auditable checks: (1) vendor‑syntax parsing into a canonical AST; (2) per-command normalization and ordering checks; (3) intent-constraint enforcement (required\_sequence and allowed\_touched\_fields); (4) dependency and transient-rule checks (e.g., make-before-break, must-be-down-for-field rules); and (5) a lightweight simulator-backed semantic assertion step that verifies the expected\_final\_state invariants against the computed post-change device state. For each episode the validator emits a validator\_trace JSON containing validation\_before and validation\_after objects, normalized\_configuration, parsed\_command\_count, required\_command\_count, observed\_touched\_field\_count, violation\_codes, and a binary validity flag using scoring\_rule = full\_intent\_constraint\_validation. Per-episode scoring JSONs and normalized responses are archived under execution/scoring/ with SHA‑256 digests (examples and digests included in the execution bundle). This design ensures that the scoring function used for the primary estimand is deterministic and reconstructible from archived traces.

Failure-treatment, retries, exclusions, and contamination detection. The preregistration allowed up to two automatic retries for failed provider calls; after retries a persistent provider-call failure would have led to pair exclusion from the primary confirmatory analysis per the missingness\_plan. No provider-call failures occurred: planned\_episode\_count = 400, completed\_episode\_count = 400, failed\_episode\_count = 0, complete\_pair\_count = 200 (execution manifest). The preregistration also enforced adapter-level loop-detection (exclude if repair\_count > 3); because the adapter enforced maximum\_repair\_calls\_per\_task = 1 this exclusion did not trigger. Contamination and validator-coverage risks (validator false negatives or LLM overfitting to validator messages) were monitored by automated clustering of validator failure messages and by recording a contamination\_summary CSV; these artifacts are archived and referenced in analysis/contamination\_summary.csv (SHA-256 available in the manifest). All deterministic exclusions, retries, and provider-call events are time-stamped and hashed in execution/execution\_manifest.json for auditability.

Primary scoring and analysis execution. The analysis pipeline used the paired\_binary\_analysis\_v1 executor. Input = execution/raw\_results.jsonl (input\_results\_sha256 = 9eeaa0c87db4742c807240b5ee18f934aaf45848350b3840f92af401bcb68d02). The executor computes the 2\ensuremath{\times}2 contingency cells (n\_11, n\_10, n\_01, n\_00) under a clear guarded-first CSV header ordering (analysis/paired\_contingency\_table.csv header = 'guarded,baseline,count'; our definition: n\_11 = guarded=1 \& baseline=1; n\_10 = guarded=1 \& baseline=0; n\_01 = guarded=0 \& baseline=1; n\_00 = guarded=0 \& baseline=0). Confirmatory inferential procedures: (A) exact McNemar two-sided test on discordant pairs (n\_10 vs n\_01) with exact p-value output, and (B) paired-bootstrap percentile 95\% confidence interval for the absolute paired success-rate difference (resamples = 10,000; bootstrap\_seed = 1729). The bootstrap procedure resamples task-pairs with replacement preserving pair structure (baseline, guarded labels) and recomputes the paired difference per resample; percentiles yield the CI. All numeric analysis outputs and logs (analysis/analysis\_log.jsonl SHA-256 = dbc0bd6b...) are archived.

Secondary and exploratory analyses (prespecified). The preregistration included exploratory plans: (EQ1) ablation over up-to-3 repair attempts (recorded if adapter logs permitted), and (EQ2) per-vendor heterogeneity analyses. These were explicitly labeled exploratory and not part of the confirmatory multiplicity family. Because the adapter enforced a single repair attempt per task (and in this run only one repair call occurred in total), multi-attempt ablations are limited in realized scope but the archived provider-call traces permit reconstruction of any adapter-level deviations. Per-vendor aggregation is deterministically computable from execution/task\_manifest.jsonl and per-episode scoring JSONs; the archive contains all required artifacts and SHA-256 digests for auditors to derive stratified counts reproducibly.

Forensic reproducibility and artifact referencing. To ensure deterministic forensic replication we publish authoritative SHA‑256 digests for principal inputs and analysis artifacts: execution/model\_configuration.json SHA-256 = a40e96e212ab87da251ac0d6dbb75bc543afa13438b5a6b9ebd073597ffdd820; execution/raw\_results.jsonl SHA-256 = 9eeaa0c87db4742c807240b5ee18f934aaf45848350b3840f92af401bcb68d02; execution/task\_manifest.jsonl SHA-256 = 4b28215e8a611feda50f6284058083b9ffcc2f5881f557081ce9e9f890f1c78a; analysis/paired\_contingency\_table.csv SHA-256 = 658051bae1f10d90aed1728b8d20dd3700a5694b68b02f1b0c118c45d9414db0; analysis/condition\_summary.csv SHA-256 = 86ab6b56e3ec10aa54aa08480a8e1ef31b64d79bf22bc99dac1d8ed00628a68a; analysis/deterministic\_reconciliation.json SHA-256 = 8a2b85f8260409eebce1641421af6a74f1c479e505bdef805314510e25931891. These exact artifact pointers enable an auditor to (a) recompute contingency cells, (b) inspect the single discordant episode validator\_trace, and (c) verify provider-call accounting that explains model\_calls\_used = 201 = 200 shared\_initial + 1 repair (execution manifest). All validator\_trace fields (validation\_before, validation\_after, violation\_codes) are recorded per episode and are sufficient to classify residual error types (semantic/policy vs syntactic/parse) deterministically.

Implementation notes and reproducibility hygiene. The pipeline components (task generator, adapter, validator, analysis executor) were implemented as modular JSON‑driven adapters. All seeds, generator ids, and versions (e.g., deterministic\_netops\_task\_generator\_v5, deterministic\_netops\_validator\_v5) are preserved in the task\_manifest and scoring JSONs. No cross-condition randomization was performed after initial shared\_initial generation. The code and JSON schemas used by the executor are archived and listed in execution/execution\_manifest.json; the execution manifest contains a full per-file hash manifest for forensic reconstruction. Any reviewer or auditor can reconstruct the exact paired inputs and outputs by fetching the archived JSON lines and validating SHA‑256 digests against the manifest.

\section{Results}
Primary confirmatory outcome (artifact-grounded). Complete\_pair\_count = 200. Baseline successes = 199/200 (baseline\_success\_rate = 0.995); guarded successes = 200/200 (guarded\_success\_rate = 1.000). Contingency counts (analysis/paired\_contingency\_table.csv SHA-256 = 658051bae1f10d90aed1728b8d20dd3700a5694b68b02f1b0c118c45d9414db0): n\_11 = 199, n\_10 = 1, n\_01 = 0, n\_00 = 0. Absolute paired change (guarded - baseline) = 0.005. Exact McNemar two-sided p = 1.0. Paired bootstrap 95\% percentile CI = [0.0, 0.015] (10,000 resamples, seed 1729) (analysis/analysis\_log.jsonl SHA-256 = dbc0bd6b...).

Estimand algebra and preregistered hypothesis reconciliation. The preregistration phrased H1 as a >=50\% relative reduction in actionable-error rate, while the archived primary\_estimand is the paired success-rate difference (guarded - baseline). Let baseline\_error = 1 - baseline\_success. A 50\% relative reduction in actionable-error rate implies absolute reduction = 0.5 \ensuremath{\times} baseline\_error. With observed baseline\_error = 0.005 (1/200), a 50\% relative reduction corresponds to an absolute change of 0.0025 --- an absolute effect smaller than our originally assumed absolute-effect target (0.15) used in power planning. Consequently, although the guarded pipeline eliminated the single observed actionable error (1\ensuremath{\rightarrow}0), the original power calibration (sized to detect large absolute effects) does not permit confirmatory affirmation of the preregistered large-effect claim. All algebraic mappings and contingency counts above are derived from archived analysis CSVs (SHA-256s in Disclosure) and are authoritative.

Representative diagnostic episode(s). The single discordant improvement (n\_10 = 1) is recorded as a unique episode in execution/raw\_results.jsonl and the per-episode scoring JSON set under execution/scoring/; the full hash manifest and episode-level traces are available for auditors (execution/execution\_manifest.json listing all per-file SHA-256 values). For forensic reproducibility auditors can deterministically locate the discordant episode and inspect its validator\_trace (validation\_before/after) using the archived raw\_results and scoring JSONs (SHA-256 digests in Disclosure). In-text we do not print the per-episode file content verbatim; instead we provide the precise artifact pointers and SHA-256 digests so auditors can retrieve the exact episode id and complete trace.

\section{Discussion}
Expanded discussion (artifact-grounded technical detail and diagnostics).

1) Reconciling estimand choice, hypothesis wording, and practical decision rules. The practical difference between a relative-error‑reduction hypothesis (">=50\% relative reduction in actionable-error rate") and the preregistered absolute paired estimand (guarded - baseline success-rate difference) is numerical and decision-relevant in low-prevalence regimes. The archived estimator used for inference and for the preregistered decision rule is the absolute paired difference (paired\_success\_rate\_difference\_guarded\_minus\_baseline), implemented and archived in analysis/paired\_contingency\_table.csv (SHA-256 = 658051bae1f10d90aed1728b8d20dd3700a5694b68b02f1b0c118c45d9414db0) and analysis/condition\_summary.csv (SHA-256 = 86ab6b56e3ec10aa54aa08480a8e1ef31b64d79bf22bc99dac1d8ed00628a68a). Analysts and designers should therefore interpret confirmatory outcomes in terms of the absolute paired difference and its CI rather than raw relative percentages when baseline\_error is small. We emphasize that the observed 1\ensuremath{\rightarrow}0 improvement yields an absolute paired change of 0.005; algebraically this corresponds to a 50\% relative reduction only if the baseline error equals 0.01; for baseline\_error=0.005 the equivalent absolute reduction to satisfy a 50\% relative reduction would be 0.0025 --- a quantity outside the scope of the preregistered power plan sized for large absolute effects. This arithmetic mapping is deterministic from the archived counts and is the correct reconciliation between the hypothesis wording and the estimator used.

2) Practical audit workflow for the discordant episode and per‑vendor aggregation. The execution repository contains every per-episode trace required to reproduce the contingency cell and to inspect the validator's before/after ASTs: execution/raw\_results.jsonl (SHA-256 = 9eeaa0c87db4742c807240b5ee18f934aaf45848350b3840f92af401bcb68d02), execution/task\_manifest.jsonl (SHA-256 = 4b28215e8a611feda50f6284058083b9ffcc2f5881f557081ce9e9f890f1c78a), and the per-episode scoring JSONs under execution/scoring/ (full hash manifest under execution/execution\_manifest.json). A reproducible forensic procedure is: (a) locate the single discordant pair in analysis/paired\_contingency\_table.csv (n\_10 = 1), (b) identify its pair\_id in execution/raw\_results.jsonl, (c) open execution/scoring/<episode>.json to inspect validation\_before and validation\_after objects (ASTs, parsed\_command\_count, violation lists), and (d) verify provider-call provenance via execution/provider\_calls/*. The repository-anchored SHA‑256 values in the Disclosure allow auditors to verify file integrity before inspection. We implemented these steps in the internal analysis executor and recorded the actions in analysis/analysis\_log.jsonl (SHA-256 = dbc0bd6b29e8624a2f72a35b0073a6ee7644daa85db4a7948a9fa75e276c6c3d). Because the artifact set is canonical, third-party auditors can deterministically reconstruct the same episode-level narrative without ambiguity.

3) Validator coverage diagnostics and residual error taxonomy. The deterministic validator produces structured validation\_before and validation\_after objects that include parsed\_command\_count, required\_command\_count, normalized\_configuration, and explicit violation\_codes. These fields enable reproducible classification of residual errors into syntactic/parse (validator cannot parse or vendor-syntax fails), workflow/sequencing (intent-ordering, make-before-break violations), and semantic/policy mismatches (simulator-detected reachability or group constraints). In this run the guarded condition produced zero residual actionable errors (guarded\_successes = 200). Because the residual count is zero, the preregistered H2 contingency test (residual error-type distribution comparison) was untestable as a confirmatory comparison; any discussion of residual error-type predominance is exploratory and should be examined on larger samples or deliberately higher baseline-error strata. The per-episode validator\_trace objects are the forensic basis for such classification and are archived per-episode under execution/scoring/ (SHA-256 pointers in Disclosure).

4) Adapter-contract semantics and stage accounting implications. The adapter enforced shared\_initial\_generation semantics (single shared candidate across conditions) and a strict repair budget (maximum\_repair\_calls\_per\_task = 1). Provider-call accounting in analysis/deterministic\_reconciliation.json (SHA-256 = 8a2b85f8260409eebce1641421af6a74f1c479e505bdef805314510e25931891) shows hosted\_model\_call\_event\_count = 201 and stage\_counts = \{"shared\_initial\_generation": 200, "repair": 1\}, matching model\_calls\_used = 201. This accounting explains why the guarded condition yielded zero additional shared initial provider-calls: the pipeline reused the same initial output and exercised the repair stage only when invoked. Designers should note that adapter budgeting interacts nonlinearly with validator coverage: with a generous repair budget multiple repair attempts could correct more baseline errors but also expose potential failure modes (LLM collusion with validator hints or overfitting to validator messages). The archived provider\_call\_audit (execution/execution\_manifest.json and deterministic\_reconciliation.json) documents each call and is the source of truth for auditing these interactions.

5) Power, estimand selection for rare-failure regimes, and design recommendations. This confirmatory run demonstrates a common practical pitfall: powering a trial for large absolute effects (e.g., \ensuremath{\Delta}\ensuremath{\approx}0.15) while the true baseline\_error is rare (here 0.005). When baseline failure prevalence is small relative to target absolute reductions, paired binary estimands produce very low event counts and wide relative uncertainty for requested relative-reduction claims. Remedies: (a) increase N; (b) oversample high-difficulty strata where baseline\_error is larger (the task generator encodes difficulty per task\_payload.difficulty and auditors can filter by it deterministically); (c) adopt weighted estimands that upweight rare catastrophic-error types or cost-weighted outcomes; or (d) pre-specify hierarchical or Bayesian estimands that borrow strength across strata. All these alternatives are compatible with the existing forensic artifact model and can be simulated deterministically from the archived task templates and generator seeds in execution/task\_manifest.jsonl.

6) Threats to construct and external validity tied to validator scope. The deterministic validator is necessarily an abstraction: it operationalizes a subset of vendor syntax, workflow policies, and semantic assertions available in the simulator. Residual unobserved semantic risks (e.g., vendor‑specific side-effects outside the validator model) remain a concern and are explicitly bounded by the synthetic task bank's policy language. Auditors should inspect validator coverage by sampling validator\_trace.violation\_codes across episodes to locate unmodeled policy types; those JSON fields are authoritative for coverage analyses.

7) Reproducibility practice and reviewer requests. Reviewers requested inline per-vendor stratified counts and the explicit discordant episode identifier. Both are deterministically derivable from the archived artifacts; the repository contains the authoritative JSONs needed to extract them (execution/task\_manifest.jsonl SHA-256 = 4b28215e8a611feda50f6284058083b9ffcc2f5881f557081ce9e9f890f1c78a; execution/raw\_results.jsonl SHA-256 = 9eeaa0c87db4742c807240b5ee18f934aaf45848350b3840f92af401bcb68d02; execution/scoring/*.json SHA-256s listed in execution/execution\_manifest.json). We preserve deterministic reproducibility by pointing reviewers to these exact artifact paths and SHA-256 digests; auditors can extract and verify per-vendor aggregates and the discordant-pair id without ambiguity. If reviewers require those numbers printed verbatim in‑manuscript we will include them in a subsequent revision only after deterministic extraction (to avoid introducing errors in reporting), citing the extracted CSV with its SHA-256 digest inline.

8) Operational implication summary. Under the conservative adapter contract and the validator implementation used, the practical marginal benefit of adding a deterministic validator is tightly coupled to baseline\_error prevalence and to the validator's semantic coverage. For organizations with very low baseline LLM error rates, validator investment should be guided by the cost of a single failure and by targeted improvements in validator semantic coverage for catastrophic categories rather than by expecting large absolute proportion reductions. For higher baseline-error contexts, deterministic validators with broader simulator-backed assertions and richer repair budgets are more likely to produce material absolute gains; these configurations can be explored reproducibly from the archived task templates and generator seeds.

In sum, the discussion expands the artifacts-to-decision pipeline: arithmetic reconciliation of estimand vs hypothesis, concrete auditing steps to locate and inspect the discordant episode and per-vendor aggregates, validator-coverage diagnostics using per-episode JSON fields, and design recommendations for future confirmatory runs in rare-failure regimes. All diagnostic steps above reference specific archived files and SHA-256 digests in the Disclosure so auditors can reproduce every assertion deterministically.

\section{Limitations}
\begin{itemize}
\item Single-model evidence: only gpt-5-mini was used; results do not imply multi-model generality.
\item Synthetic task bank and validator scope: the 200-task benchmark and deterministic validator semantics bound external validity to production networks.
\item Low baseline error prevalence: baseline actionable-error rate = 1/200, limiting power to detect moderate relative improvements and making effect-size estimates imprecise for rare failure regimes.
\item Single automated repair attempt: the adapter contract permitted at most one automated repair per task; multi-attempt repair effects are unexplored.
\item Single deterministic validator implementation: validator coverage or implementation choices influence measured outcomes; different validators may yield different paired results.
\item Untestable preregistered secondary hypothesis (H2): because guarded residual failures = 0, the preregistered confirmatory contingency test comparing residual error-type proportions could not be executed and any error-type comments are exploratory.
\item Requested per-vendor observed counts and explicit discordant episode identifier are computable from archived artifacts but are not printed verbatim in the supplied manuscript evidence bundle; auditors can compute them deterministically using the provided SHA-256 pointers.
\end{itemize}

\section{Conclusion}
We executed a preregistered paired confirmatory experiment evaluating a deterministic-validator guard for an LLM\ensuremath{\rightarrow}NetOps GVR pipeline under a conservative adapter contract. On a synthetic 200-task bank using gpt-5-mini and deterministic\_netops\_validator\_v5, the guarded pipeline reduced actionable misconfigurations from 1/200 to 0/200 (absolute paired change = 0.005). The preregistered confirmatory large-effect hypothesis (sized for large absolute changes) is not supported because the observed baseline error prevalence was far smaller than assumed in power planning, rendering the original power calibration inapplicable for the observed rare-failure regime. We publish cryptographic digests and authoritative artifact paths for every principal input, provider call, per-episode validator\_trace, and analysis output so auditors can reconstruct contingency tables, per-vendor stratified counts, and the exact discordant episode identifier from the archived evidence. Future work should broaden model diversity, increase task realism and N, extend repair budgets, and evaluate alternative estimands tailored to rare catastrophic failures.

\section*{Disclosure Statement}
Disclosure Statement (mandatory). Initial master prompt immutable reference: provenance/master\_prompt.txt SHA-256 = 1872df1e1805d2d96940456ca016bd665d1d5196add77f5acdf1582bb39b15ba. Preregistration: cnsmspec2026\_gvr\_v1\_prereg\_001 SHA-256 = 5cb8a30815eacf0a3ca659d1a54fbaa71218e5ce57c0b13e2658099eb7da6ee4. Declared model and configuration: execution/model\_configuration.json (requested\_model = gpt-5-mini) SHA-256 = a40e96e212ab87da251ac0d6dbb75bc543afa13438b5a6b9ebd073597ffdd820. Execution raw results and task manifest: execution/raw\_results.jsonl SHA-256 = 9eeaa0c87db4742c807240b5ee18f934aaf45848350b3840f92af401bcb68d02; execution/task\_manifest.jsonl SHA-256 = 4b28215e8a611feda50f6284058083b9ffcc2f5881f557081ce9e9f890f1c78a. Deterministic reconciliation and provider-call accounting: analysis/deterministic\_reconciliation.json SHA-256 = 8a2b85f8260409eebce1641421af6a74f1c479e505bdef805314510e25931891. Paired contingency evidence: analysis/paired\_contingency\_table.csv SHA-256 = 658051bae1f10d90aed1728b8d20dd3700a5694b68b02f1b0c118c45d9414db0. Condition summary (success rates): analysis/condition\_summary.csv SHA-256 = 86ab6b56e3ec10aa54aa08480a8e1ef31b64d79bf22bc99dac1d8ed00628a68a. Missingness summary: analysis/missingness\_summary.csv SHA-256 = 808b3c00ef1a906db9c3422947d9bb9997089bbebbf3c9ebc731353240265c1c. Analysis executor inputs: analysis/results.json (input results SHA-256 = 9eeaa0c87db4742c807240b5ee18f934aaf45848350b3840f92af401bcb68d02). Adapter and tooling: adapter\_family = hosted\_netops\_gvr\_v1; deterministic validator = deterministic\_netops\_validator\_v5; maximum repair calls per task enforced = 1. Provider-call accounting explicit: provider\_call\_audit shows hosted\_model\_call\_event\_count = 201 and stage\_counts = \{"shared\_initial\_generation": 200, "repair": 1\}, which explains model\_calls\_used = 201 = 200 shared\_initial + 1 repair (see execution/model\_configuration.json SHA-256 above and deterministic\_reconciliation.json SHA-256 above). Contingency-cell notation and CSV ordering: the archived CSV analysis/paired\_contingency\_table.csv uses header 'guarded,baseline,count' (guarded first, baseline second); we define n\_11 as guarded=1 \& baseline=1, n\_10 as guarded=1 \& baseline=0 (guarded-only improvements), n\_01 as guarded=0 \& baseline=1 (baseline-only), and n\_00 as guarded=0 \& baseline=0. Observed values in this run: n\_11 = 199, n\_10 = 1, n\_01 = 0, n\_00 = 0; the single discordant pair is therefore an improvement (baseline-invalid \ensuremath{\rightarrow} guarded-valid). Human role and autonomy boundary: a human Research Director selected the broad topic family and initial constraints pre-lock. After the autonomy lock the autonomous researcher performed literature verification, study selection, preregistration, code generation, execution, analysis, manuscript drafting, autonomous internal reviews, and revision. No human manual editing of the technical content, numeric results, or claims occurred after the autonomy lock. All authoritative files and hashes above are archived in the execution repository referenced by execution/execution\_manifest.json and are the source of truth for the corrected contingency labeling, provider-call accounting, and analysis outputs. Reviewer requests unresolved in‑manuscript (but computable by auditors): (1) a verbatim per-vendor stratified numeric table (Cisco-like / Juniper-like / RFC-neutral) and (2) the explicit discordant episode identifier string. Both values are derivable from execution/task\_manifest.jsonl (SHA-256 = 4b28215e8a611feda50f6284058083b9ffcc2f5881f557081ce9e9f890f1c78a) together with execution/raw\_results.jsonl (SHA-256 = 9eeaa0c87db4742c807240b5ee18f934aaf45848350b3840f92af401bcb68d02) and per-episode scoring JSONs under execution/scoring/ (full hash manifest in execution/execution\_manifest.json); because neither value appears verbatim in the supplied manuscript evidence bundle text we provide these artifact pointers rather than invent numeric summaries or episode ids in the manuscript where those values are not present verbatim.

\begin{thebibliography}{99}
\bibitem{ref1} http://citeseerx.ist.psu.edu/viewdoc/summary?doi=10.1.1.228.5064
\bibitem{ref2} https://doi.org/10.1145/2342441.2342452
\bibitem{ref3} https://doi.org/10.1145/3656296
\bibitem{ref4} https://doi.org/10.1109/mcom.001.2400368
\end{thebibliography}

\end{document}
